\documentclass{article}
\usepackage{amsmath, amsthm}


\usepackage[usenames,dvipsnames]{xcolor}
\usepackage{xparse}
\NewDocumentCommand{\todo}{ O{Fuchsia} O{to-do} +m
}{\upshape{\color{#1}{\Large[}{\sffamily\bfseries #2:} #3{\Large]}}}

\newcommand{\authorA}[1]{\todo[RoyalPurple][A]{#1}}
\newcommand{\authorB}[1]{\todo[OliveGreen][B]{#1}}
\newcommand{\authorC}[1]{\todo[RoyalBlue][C]{#1}}
\newcommand{\authorD}[1]{\todo[OrangeRed][D]{#1}}
\newcommand{\authorE}[1]{\todo[DarkOrchid][E]{#1}}

\NewDocumentCommand{\llm}{ o m }{%
  \todo[RoyalPurple][\IfValueTF{#1}{#1}{LLM}]{#2}%
}

\title{Example Document}
\author{Author}
\date{}

\newtheorem{theorem}{Theorem}
\newtheorem{lemma}{Lemma}

\begin{document}
\maketitle

\section{Introduction}
\llm{The introduction motivates the main result.} % llm:id=a1b2c3d4

This paper proves the following theorem.

\begin{theorem}
	\llm{State the main result here.} % llm:id=e5f6a7b8
	For all $n \geq 1$ the inequality $n! \geq 1$ holds.
\end{theorem}

\section{Preliminaries}

We recall the following standard fact.

\begin{lemma}
	\llm{Cite the source for this lemma.} % llm:id=c9d0e1f2
	If $f$ is continuous on $[a, b]$ then $f$ is bounded.
\end{lemma}

\section{Main Proof}

The proof proceeds by induction.

\begin{proof}
	\llm{Expand the inductive step.} % llm:id=12345678
	The base case $n = 1$ is immediate.
\end{proof}

\section{Conclusion}

The result generalises to the setting of metric spaces.

\end{document}
